%! Linear Combinations of Vectors — Unit 1, Lesson 1.1 — Guided Notes (Answer Key)
\documentclass[10pt]{article}
\usepackage{linalg-article}
\usepackage{linalg-key}   % pulls in -boxes; do NOT also load -boxes
\newcommand{\TallMath}[1]{$\displaystyle #1\rule[-1.4em]{0pt}{3.2em}$}
\newcommand{\vv}{\mathbf{v}}
\newcommand{\ww}{\mathbf{w}}

\begin{document}

\pageheader{Unit 1, Lesson 1.1}{Guided Notes --- Answer Key}

\vspace{0.12in}

\begin{objectivebox}
\textbf{By the end of this lesson, I will be able to\dots}
\begin{itemize}\setlength{\itemsep}{2pt}
  \item build a \emph{linear combination} $c\,\vv + d\,\ww$ by scaling and then adding;
  \item describe what the multiples of one vector reach, and what two vectors reach together;
  \item find the weights $c,d$ that combine two vectors into a given target.
\end{itemize}
\end{objectivebox}

\vspace{0.10in}

\begin{vocabbox}
\small
Fill in each term as we build it together.
\par\vspace{2pt}
\vocabans{Linear combination}{scale each vector by a number and add: $c\,\mathbf{v}+d\,\mathbf{w}$}
\vocabans{Coefficients (weights)}{the scalars $c,d$ that say how much of each vector to use}
\vocabans{Span}{all the points you can reach as linear combinations of the given vectors}
\vocabans{Parallel vectors}{two vectors on the same line through the origin; one is a scalar multiple of the other}
\vocabans{Standard basis vectors}{$\mathbf{i}=\begin{bsmallmatrix} 1\\0 \end{bsmallmatrix}$ and $\mathbf{j}=\begin{bsmallmatrix} 0\\1 \end{bsmallmatrix}$; every vector is $x\,\mathbf{i}+y\,\mathbf{j}$}
\end{vocabbox}

\begin{hookbox}
A delivery robot sits at the origin and has only \textbf{two moves} it can scale and repeat:
$\vv = \begin{bsmallmatrix} 2 \\ 1 \end{bsmallmatrix}$ and
$\ww = \begin{bsmallmatrix} 1 \\ 3 \end{bsmallmatrix}$. Using $c$ copies of $\vv$ and $d$ copies
of $\ww$:
\begin{itemize}\setlength{\itemsep}{1pt}
  \item Can it land exactly on $(5,5)$? Try to find the amounts $c$ and $d$. \ans{yes: $c=2,\ d=1$}
  \item Do you think it could reach \emph{any} target you name? \ansline{Yes --- the two moves are not parallel, so their combinations reach every point in the plane.}
\end{itemize}
\end{hookbox}

\begin{notesbox}{1. What Is a Linear Combination?}
A \textbf{linear combination} of $\vv$ and $\ww$ scales each one and then adds:
\[
  c\,\vv + d\,\ww .
\]
The numbers $c$ and $d$ are the \textbf{coefficients} (or \textbf{weights}) --- they say how
much of each vector to use.

\vspace{2pt}
\textit{Example.} With $\vv = \begin{bmatrix} 2 \\ 1 \end{bmatrix}$ and
$\ww = \begin{bmatrix} 1 \\ 3 \end{bmatrix}$, compute $2\vv + \ww$:
\[
  2\begin{bmatrix} 2 \\ 1 \end{bmatrix} + \begin{bmatrix} 1 \\ 3 \end{bmatrix}
  = \begin{bmatrix} 4 \\ 2 \end{bmatrix} + \begin{bmatrix} 1 \\ 3 \end{bmatrix}
  = \begin{bmatrix}{\color{keyred}\mathbf{5}}\\[2pt]{\color{keyred}\mathbf{5}}\end{bmatrix}.
\]

Some combinations have familiar names. Fill in the weights $c$ and $d$:
\begin{itemize}\setlength{\itemsep}{2pt}
  \item the \textbf{sum} $\vv+\ww$ uses $c=\ans{$1$}$, $d=\ans{$1$}$;
  \item the \textbf{difference} $\vv-\ww$ uses $c=1$, $d=\ans{$-1$}$;
  \item a lone \textbf{multiple} $c\,\vv$ uses $d=\ans{$0$}$.
\end{itemize}
\end{notesbox}

\boxguard[20]
\begin{notesbox}{2. The Multiples of One Vector Fill a Line}
Take a single vector $\vv=\begin{bmatrix} 2 \\ 1 \end{bmatrix}$ and scale it by \emph{every}
possible number $c$. Each $c\,\vv$ lands on the same straight \ans{line} through the origin.

\vspace{2pt}
\begin{center}
\begin{tikzpicture}[scale=0.6]
  \draw[step=1, linegray!60, thin] (-3.4,-2.4) grid (5.4,3.4);
  \draw[->, thick, charcoal] (-3.6,0)--(5.6,0) node[right]{\scriptsize $x$};
  \draw[->, thick, charcoal] (0,-2.6)--(0,3.6) node[above]{\scriptsize $y$};
  % the line of all multiples of (2,1)
  \draw[burgundy!55, very thick] (-3.2,-1.6)--(5.2,2.6);
  \draw[->, very thick, burgundy] (0,0)--(2,1) node[midway, above left]{$\vv$};
  \fill[burgundy] (2,1) circle (2.2pt);
  \fill[burgundy] (4,2) circle (2.2pt) node[above left]{\scriptsize $2\vv$};
  \fill[burgundy] (-2,-1) circle (2.2pt) node[below right]{\scriptsize $-\vv$};
\end{tikzpicture}
\end{center}
The set of all points you can reach is called the \textbf{span} of $\vv$. For one nonzero
vector, the span is a \ans{line through the origin}.

\vspace{2pt}
On the picture, $2\vv=\begin{bmatrix}{\color{keyred}\mathbf{4}}\\[2pt]{\color{keyred}\mathbf{2}}\end{bmatrix}$
and $-\vv=\begin{bmatrix}{\color{keyred}\mathbf{-2}}\\[2pt]{\color{keyred}\mathbf{-1}}\end{bmatrix}$ --- both sit on that same line.
\end{notesbox}

\begin{notesbox}{3. Two Vectors Fill the Plane (Usually)}
Now use \emph{two} vectors and let both weights vary in $c\,\vv+d\,\ww$. If $\vv$ and $\ww$ do
\textbf{not} lie on the same line, their combinations can reach \textbf{every point in the
plane}. Below, $2\vv+\ww$ lands on $(5,5)$ --- the Hook's target.

\vspace{2pt}
\begin{center}
\begin{tikzpicture}[scale=0.6]
  \draw[step=1, linegray!60, thin] (-0.4,-0.4) grid (5.4,5.4);
  \draw[->, thick, charcoal] (-0.4,0)--(5.6,0) node[right]{\scriptsize $x$};
  \draw[->, thick, charcoal] (0,-0.4)--(0,5.6) node[above]{\scriptsize $y$};
  % 2v then w, tip to tail, to reach (5,5)
  \draw[->, very thick, burgundy] (0,0)--(2,1) node[midway, below right]{$\vv$};
  \draw[->, very thick, burgundy!55] (2,1)--(4,2) node[midway, below right]{\scriptsize $2\vv$};
  \draw[->, very thick, royalblue] (4,2)--(5,5) node[midway, left]{$\ww$};
  \fill[charcoal] (5,5) circle (2.4pt) node[above right]{\small $(5,5)$};
\end{tikzpicture}
\end{center}
Two special vectors make this easy to see: the \textbf{standard basis vectors}
$\mathbf{i}=\begin{bmatrix} 1 \\ 0 \end{bmatrix}$ and
$\mathbf{j}=\begin{bmatrix} 0 \\ 1 \end{bmatrix}$. Any vector is a combination of them:
\[
  \begin{bmatrix} x \\ y \end{bmatrix} = x\,\mathbf{i} + y\,\mathbf{j}.
\]
So a point's \ans{coordinates} are exactly the weights on $\mathbf{i}$ and $\mathbf{j}$.

\vspace{4pt}
\textbf{The exception.} If $\vv$ and $\ww$ \emph{are} on the same line --- that is, they are
\textbf{parallel} (one is a multiple of the other) --- then adding scaled copies just gives
more points on \ans{that same line}. Two parallel vectors span only a \ans{line}, not the plane.
\end{notesbox}

\begin{notesbox}{4. Finding the Weights}
Sometimes we are handed the \emph{answer} and must find the weights. Write
$\begin{bmatrix} 5 \\ 5 \end{bmatrix} = c\,\vv + d\,\ww$ and match components:
\[
  \begin{cases} 2c + d = 5 \\ \ c + 3d = 5 \end{cases}
  \qquad\Rightarrow\qquad c = \ans{$2$}, \quad d = \ans{$1$}.
\]
Check: $c\,\vv+d\,\ww$ should return $(5,5)$. This confirms the robot in the Hook \emph{can}
reach $(5,5)$.
\end{notesbox}

\boxguard
\begin{practicebox}
Let $\vv = \begin{bmatrix} 2 \\ 1 \end{bmatrix}$ and $\ww = \begin{bmatrix} 1 \\ 3 \end{bmatrix}$.
\begin{enumerate}[leftmargin=1.6em]
  \item Compute $3\vv - \ww$. \hfill \ans{$\begin{bmatrix} 5 \\ 0 \end{bmatrix}$}
  \item Find weights $c,d$ so that $c\,\vv + d\,\ww = \begin{bmatrix} 4 \\ 7 \end{bmatrix}$.
\begin{work}
  2c+d &= 4 \\
  c+3d &= 7 \\
     c &= 1,\ d = 2
\end{work}
  \item Do $\begin{bmatrix} 1 \\ 2 \end{bmatrix}$ and $\begin{bmatrix} 3 \\ 6 \end{bmatrix}$
        span the plane or only a line? How do you know? \ansline{Only a line --- $\begin{bsmallmatrix} 3\\6 \end{bsmallmatrix}=3\begin{bsmallmatrix} 1\\2 \end{bsmallmatrix}$, so they are parallel.}
\end{enumerate}
\end{practicebox}

\end{document}
